\clearpage\section{MOS transistor schematic}
These editable vector schematics connect equal node names within a circuit. Primitive-cell names are local pins; the hierarchy defines their connections. B denotes the body terminal; 0 is global ground. Each panel identifies its netlist and scope; measurement sources are connected as described.
Exact cells and hierarchy for \texttt{01\_mux\_truth\_table.sp} and \texttt{02\_mux\_delay\_load.sp}.
\par
\begin{figure}[!ht]
\centering
\begin{circuitikz}[font=\small]
\node[pmos,bulk,rotate=-90] (TP) at (0,1) {};
\node[nmos,bulk,rotate=90] (TN) at (0,-1) {};
\draw (TP.D) -- (TP.D -| -2,0) -- (TN.D -| -2,0) -- (TN.D);
\draw (TP.S) -- (TP.S -| 2,0) -- (TN.S -| 2,0) -- (TN.S);
\draw (-2,0) node[circ]{} -- (-2.6,0) node[left]{a (D)};
\draw (2,0) node[circ]{} -- (2.6,0) node[right]{z (S)};
\draw (TP.G) -- ++(0,.5) node[above]{eb};
\draw (TN.G) -- ++(0,-.5) node[below]{en};
\draw (TP.bulk) -- ++(0,-.35) -- ++(.8,0) node[right,font=\scriptsize]{B=vp};
\draw (TN.bulk) -- ++(0,.35) -- ++(.8,0) node[right,font=\scriptsize]{B=vn};
\node[font=\scriptsize] at (-.8,1.4){TP};
\node[font=\scriptsize] at (-.8,-1.4){TN};
\end{circuitikz}
\caption{Pass-gate cell used in both TG architectures.}
\end{figure}
\par
\par
\begin{figure}[!ht]
\centering
\begin{circuitikz}[font=\small]
\node[pmos,bulk] (MP) at (0,2.5) {};
\node[anchor=west,font=\scriptsize] at ($(MP.center)+(.55,.40)$) {MP};
\draw (MP.bulk) -- ++(.30,0) node[right,font=\scriptsize] {B=vp};
\node[nmos,bulk] (MN) at (0,0.2) {};
\node[anchor=west,font=\scriptsize] at ($(MN.center)+(.55,.40)$) {MN};
\draw (MN.bulk) -- ++(.30,0) node[right,font=\scriptsize] {B=vn};
\draw (MP.G) -- (-1,2.5) -- (-1,.2) -- (MN.G); \draw (-1,1.35) node[circ]{} -- (-1.6,1.35) node[left]{a}; \draw (MP.D |- 0,1.35) node[circ]{} -- (2,1.35) node[right]{z};
\draw (MP.S) -- ++(0,0.5) node[above,font=\small] {\texttt{vp}};
\draw (MP.D) -- (MN.D) ;
\draw (MN.S) -- ++(0,-0.5) node[below,font=\small] {\texttt{vn}};
\end{circuitikz}
\caption{Input/output and local select-complement inverter.}
\end{figure}
\par
The inv cell has $W_P=3.6k\,\mu$m, $W_N=1.8k\,\mu$m and $L=0.25\,\mu$m. The tg cell uses the same widths/length. For tg, both device drains are a and both sources are z; the pass path is bidirectional. TP/TN correspond to MP/MN inside the tg subcircuit. PMOS bodies remain vp and NMOS bodies vn, independent of the signal terminals. The n-channel gate is en, and the p-channel gate is eb. Equal node labels connect the two pass devices in parallel.
\clearpage
\par
\begin{figure}[!ht]
\centering
\begin{circuitikz}[font=\small]
\node[draw,rounded corners=2pt,align=center,minimum width=1.6cm,minimum height=.9cm] (T0) at (1,3) {XFT0\\TG: fb0};
\draw (-1,3) node[left]{f\_in0} -- (T0.west);
\node[draw,rounded corners=2pt,align=center,minimum width=1.6cm,minimum height=.9cm] (T1) at (1,1) {XFT1\\TG: s0};
\draw (-1,1) node[left]{f\_in1} -- (T1.west);
\node[draw,rounded corners=2pt,align=center,minimum width=1.6cm,minimum height=.9cm] (T2) at (1,-1) {XFT2\\TG: fb0};
\draw (-1,-1) node[left]{f\_in2} -- (T2.west);
\node[draw,rounded corners=2pt,align=center,minimum width=1.6cm,minimum height=.9cm] (T3) at (1,-3) {XFT3\\TG: s0};
\draw (-1,-3) node[left]{f\_in3} -- (T3.west);
\node[draw,rounded corners=2pt,align=center,minimum width=1.6cm,minimum height=.9cm] (T4) at (5,2) {XFT4\\TG: fb1};
\node[draw,rounded corners=2pt,align=center,minimum width=1.6cm,minimum height=.9cm] (T5) at (5,-2) {XFT5\\TG: s1};
\draw (T0.east) -- (3,3) -- (3,2) node[circ,label=above right:fn0]{} -- (T4.west); \draw (T1.east) -| (3,2); \draw (T2.east) -- (3,-1) -- (3,-2) node[circ,label=above right:fn1]{} -- (T5.west); \draw (T3.east) -| (3,-2); \draw (T4.east) -- (7,2) -- (7,0) node[circ]{} -- (8,0) node[right]{of}; \draw (T5.east) -| (7,0);
\end{circuitikz}
\caption{Optimized TG tree: the label after TG is en, eb is its local complement. All six cells use k=8.}
\end{figure}
\par
fb0/fb1 are generated by two inv cells from s0/s1 using vf/0 rails. Every TG in this tree has body rails vf/0. CLf connects of to 0. No output inverter is present: six TGs plus two select inverters give 16 MOS. The input probes VDfi connect f\_ini (positive) to ai (negative) at zero volts. The corresponding g/c input probes use the same convention.
\clearpage
\par
\begin{figure}[!ht]
\centering
\begin{circuitikz}[font=\small]
\node[draw,rounded corners=2pt,align=center,minimum width=1.8cm,minimum height=.9cm] (A) at (0,0) {g\_ini\\$i=0\ldots3$};
\node[draw,rounded corners=2pt,align=center,minimum width=1.8cm,minimum height=.9cm] (B) at (3,0) {XGIi: inv};
\node[draw,rounded corners=2pt,align=center,minimum width=1.8cm,minimum height=.9cm] (C) at (6,0) {XGTia\\TG: $S_{1i}$};
\node[draw,rounded corners=2pt,align=center,minimum width=1.8cm,minimum height=.9cm] (D) at (9,0) {XGTib\\TG: $S_{0i}$};
\draw[-{Latex}] (A.east) -- (B.west);
\draw[-{Latex}] (B.east) -- (C.west) node[midway,above]{gai};
\draw[-{Latex}] (C.east) -- (D.west) node[midway,above]{gni};
\draw (D.east) -- (11,0) node[right]{gz};
\node[draw,rounded corners=2pt,align=center,minimum width=1.2cm,minimum height=.9cm] (E) at (3,-2) {gz};
\node[draw,rounded corners=2pt,align=center,minimum width=1.8cm,minimum height=.9cm] (F) at (6,-2) {XGO: inv};
\node[draw,rounded corners=2pt,align=center,minimum width=1.2cm,minimum height=.9cm] (G) at (9,-2) {og};
\draw[-{Latex}] (E.east) -- (F.west);
\draw[-{Latex}] (F.east) -- (G.west);
\end{circuitikz}
\caption{Lecture TG architecture: four identical indexed branches join at gz; a single output inverter regenerates og.}
\end{figure}
\par
$S_{1i}=s1$ for $i\ge2$, otherwise gb1; $S_{0i}=s0$ for odd i, otherwise gb0. eb is the opposite select literal. All cells use k=1 and vg/0 rails; two local select inverters are included. The four branch inputs and output each have an inverter, giving 30 MOS including select inversion.
\par
\begin{figure}[!ht]
\centering
\begin{circuitikz}[font=\small]
\node[draw,rounded corners=2pt,align=center,minimum width=3.5cm,minimum height=.9cm] (N) at (0,0) {$(S_{1i},S_{0i},c\_ini)$};
\node[draw,rounded corners=2pt,align=center,minimum width=2.2cm,minimum height=.9cm] (P) at (4,0) {XCN$i$: NAND3\\$i=0\ldots3$};
\node[draw,rounded corners=2pt,align=center,minimum width=3.5cm,minimum height=.9cm] (Q) at (8,0) {XCO: NAND4\\$(cn0,cn1,cn2,cn3)$};
\draw[-{Latex}] (N.east) -- (P.west);
\draw[-{Latex}] (P.east) -- (Q.west) node[midway,above]{cni};
\draw (Q.east) -- (10.5,0) node[right]{oc};
\end{circuitikz}
\caption{Static CMOS MUX hierarchy; all four NAND3 outputs enter one NAND4.}
\end{figure}
\par
Here the select complements are cb0/cb1, generated locally with vc/0 rails. All static gates use k=1. Outputs og/oc/of each have the same cl load to 0.
\clearpage
\par
\begin{figure}[!ht]
\centering
\begin{circuitikz}[font=\small]
\node[pmos,bulk] (MP0) at (0.0,3) {};
\node[anchor=west,font=\scriptsize] at ($(MP0.center)+(.55,.40)$) {MP0};
\draw (MP0.bulk) -- ++(.30,0) node[right,font=\scriptsize] {B=vp};
\draw (MP0.G) -- ++(-.45,0) node[left,font=\small] {\texttt{a0}};
\draw (MP0.S) -- (MP0.S |- 0,4.2); \draw (MP0.D) -- (MP0.D |- 0,1.7);
\node[nmos,bulk] (MN0) at (3.8,0.0) {};
\node[anchor=west,font=\scriptsize] at ($(MN0.center)+(.55,.40)$) {MN0};
\draw (MN0.bulk) -- ++(.30,0) node[right,font=\scriptsize] {B=vn};
\draw (MN0.G) -- ++(-.45,0) node[left,font=\small] {\texttt{a0}};
\node[pmos,bulk] (MP1) at (3.8,3) {};
\node[anchor=west,font=\scriptsize] at ($(MP1.center)+(.55,.40)$) {MP1};
\draw (MP1.bulk) -- ++(.30,0) node[right,font=\scriptsize] {B=vp};
\draw (MP1.G) -- ++(-.45,0) node[left,font=\small] {\texttt{a1}};
\draw (MP1.S) -- (MP1.S |- 0,4.2); \draw (MP1.D) -- (MP1.D |- 0,1.7);
\node[nmos,bulk] (MN1) at (3.8,-2.2) {};
\node[anchor=west,font=\scriptsize] at ($(MN1.center)+(.55,.40)$) {MN1};
\draw (MN1.bulk) -- ++(.30,0) node[right,font=\scriptsize] {B=vn};
\draw (MN1.G) -- ++(-.45,0) node[left,font=\small] {\texttt{a1}};
\draw (MN0.S) -- (MN1.D) node[midway,above right,font=\small] {\texttt{n1}};
\node[pmos,bulk] (MP2) at (7.6,3) {};
\node[anchor=west,font=\scriptsize] at ($(MP2.center)+(.55,.40)$) {MP2};
\draw (MP2.bulk) -- ++(.30,0) node[right,font=\scriptsize] {B=vp};
\draw (MP2.G) -- ++(-.45,0) node[left,font=\small] {\texttt{a2}};
\draw (MP2.S) -- (MP2.S |- 0,4.2); \draw (MP2.D) -- (MP2.D |- 0,1.7);
\node[nmos,bulk] (MN2) at (3.8,-4.4) {};
\node[anchor=west,font=\scriptsize] at ($(MN2.center)+(.55,.40)$) {MN2};
\draw (MN2.bulk) -- ++(.30,0) node[right,font=\scriptsize] {B=vn};
\draw (MN2.G) -- ++(-.45,0) node[left,font=\small] {\texttt{a2}};
\draw (MN1.S) -- (MN2.D) node[midway,above right,font=\small] {\texttt{n2}};
\draw (MP0.S |- 0,4.2) -- (MP2.S |- 0,4.2) node[midway,above]{vp}; \draw (MP0.D |- 0,1.7) -- (MP2.D |- 0,1.7); \draw (MN0.D) -- (MN0.D |- 0,1.7) node[circ,label=above right:z]{};
\draw (MN2.S) -- ++(0,-0.4) node[below,font=\small] {\texttt{vn}};
\end{circuitikz}
\caption{Exact NAND3 transistor cell: three parallel PMOS branches and the n1/n2 NMOS series stack.}
\end{figure}
\par
For nand3, $W_P=3.6k\,\mu$m and $W_N=5.4k\,\mu$m, $L=0.25\,\mu$m. nand4 extends the same parallel-PMOS/series-NMOS pattern to four inputs, with $W_N=7.2k\,\mu$m and the same PMOS width and length. All PMOS bodies/sources connect to vp; every NMOS body connects to vn, while intermediate sources are the stack nodes. Thus four NAND3 cells (24 MOS), one NAND4 (8 MOS), and two select inverters (4 MOS) give 36 MOS.
